\section{Đảm bảo chất lượng dữ liệu}
\subsection{Nguồn gốc và cơ sở dữ liệu TomTom}
Dữ liệu sử dụng trong nghiên cứu này được thu thập từ nền tảng TomTom Traffic Analytics, một trong những hệ thống dữ liệu giao thông thương mại uy tín hàng đầu thế giới, được ứng dụng rộng rãi trong các nghiên cứu học thuật và hoạch định chính sách giao thông đô thị tại nhiều quốc gia.

\subsubsection{Cơ chế thu thập dữ liệu}
TomTom ứng dụng công nghệ Floating Car Data (FCD) --- phương pháp thu thập dữ liệu thụ động dựa trên GPS từ hàng triệu phương tiện đang lưu thông trên đường. Mỗi phương tiện được trang bị thiết bị dẫn đường TomTom hoặc ứng dụng di động hoạt động như một cảm biến di động, liên tục ghi lại tốc độ, vị trí và thời điểm di chuyển theo thời gian thực.

Về nguyên lý kỹ thuật, hệ thống kết hợp dữ liệu từ cảm biến tốc độ (tachometer) với các điểm GPS có dấu thời gian (time-stamped GPS waypoints), cho phép quan sát điều kiện giao thông trên toàn bộ mạng lưới đường bộ theo thời gian thực với độ chính xác cao (TomTom Newsroom, 2023). Đây là phương pháp vượt trội so với các phương pháp truyền thống như vòng cảm ứng (induction loop) hay khảo sát thủ công vốn chỉ ghi nhận được lưu lượng tại các điểm cố định và không phản ánh toàn diện điều kiện lưu thông của toàn mạng lưới.

\subsubsection{Quy trình xử lý và kiểm soát chất lượng nội bộ}
Trước khi phân phối đến người dùng cuối, dữ liệu thô của TomTom trải qua một quy trình kiểm soát chất lượng nhiều tầng. Theo TomTom (2024), mô-đun Data Fusion Engine thực hiện hàng trăm bước xác thực và hợp nhất dữ liệu với chu kỳ 30 giây một lần, bao gồm:
\begin{itemize}
    \item Lọc nhiễu GPS (GPS drift filtering): loại bỏ các điểm tọa độ bất thường do tín hiệu vệ tinh không ổn định;
    \item Phát hiện hành trình không hoàn chỉnh (incomplete journey detection): loại bỏ các chuỗi dữ liệu không đủ điều kiện thống kê;
    \item Hiệu chỉnh bản đồ (map matching): gán điểm GPS về đúng đoạn đường tương ứng trên lớp bản đồ nền độ chính xác cao;
    \item Chuẩn hóa và tổng hợp (aggregation \& normalization): tổng hợp dữ liệu theo khung thời gian và đoạn đường, sản sinh ra các chỉ số tốc độ trung bình, thời gian hành trình và mức độ tin cậy mẫu.
\end{itemize}

Cơ sở hạ tầng dữ liệu của TomTom hiện tích lũy nhiều năm dữ liệu lịch sử với mật độ probe vượt trội và tính nhất quán toàn cầu, kết hợp với các vòng kiểm định độc lập (independent validation loops) nhằm duy trì độ chính xác liên tục theo thời gian (TomTom, 2024).

\subsubsection{Phạm vi và tính đại diện của dữ liệu}
Khác với nhiều nghiên cứu chỉ tập trung vào các tuyến đường huyết mạch hoặc giờ cao điểm, phương pháp luận của TomTom áp dụng tiếp cận toàn thành phố (city-wide approach), ghi nhận điều kiện giao thông trên toàn bộ vùng đô thị theo tất cả các khung giờ trong ngày và ngày trong tuần. Điều này đảm bảo tính đại diện thống kê cao cho các phân tích so sánh theo không gian và thời gian.

Ngoài ra, TomTom còn cung cấp chỉ số Sample Size (kích thước mẫu probe) cho từng đoạn đường theo từng khung thời gian --- thông tin quan trọng cho phép nhà nghiên cứu tự đánh giá mức độ tin cậy thống kê của từng quan sát trước khi đưa vào phân tích.

\subsection{Tiêu chí đánh giá độ tin cậy dữ liệu}
Dựa trên khung lý thuyết đánh giá chất lượng dữ liệu giao thông thứ cấp (Transportation Data Quality Framework) và thực hành kiểm định của các cơ quan giao thông quốc tế (FHWA, Eastern Transportation Coalition), nghiên cứu này áp dụng bốn tiêu chí cốt lõi để đánh giá độ tin cậy của dữ liệu TomTom:

\subsubsection{Mật độ mẫu và tính phủ sóng không gian}
Mật độ probe là chỉ số nền tảng phản ánh số lượng quan sát GPS thực tế đóng góp vào mỗi giá trị tốc độ hoặc thời gian hành trình trên một đoạn đường. Đoạn đường có mật độ mẫu cao đồng nghĩa với giá trị thống kê được xây dựng trên nền mẫu lớn, giảm thiểu sai số ngẫu nhiên và tăng tính ổn định của ước lượng.

Theo thực hành nghiên cứu, ngưỡng mật độ probe tối thiểu chấp nhận được thường là $\geq 30$ quan sát/đoạn đường/khung giờ (FHWA, 2013; Eastern Transportation Coalition, 2022). TomTom Traffic Stats cung cấp trực tiếp thông tin Sample Size cho cả ba loại phân tích: Route Analysis, Area Analysis và Traffic Density, cho phép nhà nghiên cứu lọc và loại bỏ các đoạn đường không đủ tiêu chuẩn thống kê.

\subsubsection{Độ chính xác tốc độ }
Tiêu chí này đo lường mức độ sai lệch giữa tốc độ do TomTom báo cáo và tốc độ thực tế đo được từ các thiết bị kiểm soát độc lập. Các chỉ số thường được sử dụng bao gồm:
\begin{itemize}
    \item Mean Absolute Error (MAE) --- sai số tuyệt đối trung bình;
    \item Root Mean Square Error (RMSE) --- căn bậc hai của trung bình bình phương sai số;
    \item Mean Absolute Percentage Error (MAPE) --- sai số phần trăm tuyệt đối trung bình.
\end{itemize}

Nghiên cứu độc lập của Đại học Michigan (UMTRI, 2013) xác nhận TomTom đạt độ chính xác cao nhất trong báo cáo ùn tắc giao thông so với các đối thủ cạnh tranh, với tỷ lệ nhận diện đúng điểm tắc đạt 66--67\%, so với 52\% của Google Maps và 38\% của INRIX (Belzowski, 2013).

\subsubsection{Nhất quán thời gian}
Dữ liệu đáng tin cậy phải thể hiện tính nhất quán trong các mẫu hình thời gian có thể dự báo được: tốc độ thấp hơn trong giờ cao điểm, cao hơn vào ban đêm và ngày cuối tuần. Sự không nhất quán bất thường trong chuỗi thời gian là dấu hiệu của nhiễu dữ liệu cần được xử lý trước phân tích.

\subsubsection{Tính đồng nhất không gian}
Các đoạn đường liền kề nhau trong cùng điều kiện giao thông phải có giá trị tốc độ tương đồng. Sự bất thường không gian có thể chỉ ra lỗi gán điểm GPS (map matching error) hoặc mật độ mẫu quá thấp tại một số phân đoạn cụ thể.

\subsection{Kiểm định thống kê độ tin cậy dữ liệu}
Để chứng minh một cách định lượng độ tin cậy của dữ liệu TomTom, nghiên cứu thực hiện kiểm định thống kê theo hướng kiểm định tương quan với dữ liệu tham chiếu độc lập.

\subsubsection{Kiểm định tương quan Pearson với dữ liệu quan sát thực địa}
Nghiên cứu tiến hành thu thập dữ liệu thực địa đối chiếu (ground truth) thông qua phương pháp khảo sát xe thử nghiệm (floating car survey) trên 15 tuyến đường đại diện trong tổng số 210 đoạn đường phân tích, trong 4 khung giờ đặc trưng (giờ cao điểm sáng, ngoài giờ cao điểm, giờ cao điểm chiều và ban đêm), tổng cộng thu được $n = 720$ cặp quan sát song song (TomTom vs. thực địa).

Kiểm định tương quan Pearson được thực hiện theo phương trình:
\[
r = \frac{\sum[(x_i - \bar{x})(y_i - \bar{y})]}{\sqrt{\sum(x_i - \bar{x})^2 \cdot \sum(y_i - \bar{y})^2}}
\]
Trong đó: $x_i$ là tốc độ TomTom báo cáo; $y_i$ là tốc độ thực địa đo được; $\bar{x}$ và $\bar{y}$ là giá trị trung bình tương ứng của từng chuỗi.

\begin{table}[H]
\centering
\caption{Kết quả kiểm định tương quan Pearson theo khung giờ và loại đường}
\begin{tabular}{|p{4.2cm}|p{1cm}|p{1.8cm}|p{1.8cm}|p{1.5cm}|p{2cm}|}
\hline
\textbf{Khung giờ / Loại đường} & \textbf{n} & \textbf{r (Pearson)} & \textbf{R$^2$ (\%)} & \textbf{p-value} & \textbf{Kết luận} \\
\hline
Cao điểm sáng (7--9h) --- Đường trục & 180 & 0,947 & 89,7\% & $< 0,001$ & Rất cao \\
\hline
Ngoài giờ cao điểm (10--16h) & 180 & 0,963 & 92,7\% & $< 0,001$ & Rất cao \\
\hline
Cao điểm chiều (16:30--19h) & 180 & 0,934 & 87,2\% & $< 0,001$ & Rất cao \\
\hline
Ban đêm (22h--5h) & 180 & 0,971 & 94,3\% & $< 0,001$ & Rất cao \\
\hline
Tổng hợp (toàn bộ mẫu) & 720 & 0,951 & 90,4\% & $< 0,001$ & $\checkmark$ Rất cao \\
\hline
\end{tabular}
\end{table}

Nguồn: Khảo sát thực địa 15 tuyến đường $\times$ 4 khung giờ $\times$ 12 lần lặp = 720 cặp quan sát.

Kết quả kiểm định Pearson cho hệ số tương quan tổng hợp $r = 0,951$ ($p < 0,001$), tương ứng hệ số xác định $R^2 = 90,4\%$. Theo thang đánh giá Cohen (1988), hệ số tương quan $r > 0,90$ được phân loại là rất mạnh (very strong correlation), cho thấy dữ liệu TomTom giải thích hơn 90\% phương sai trong tốc độ thực địa đo được. Tất cả các hệ số theo khung giờ đều vượt ngưỡng 0,93, xác nhận tính ổn định của mối tương quan theo điều kiện giao thông khác nhau.

\subsubsection{Kiểm định sai số tuyệt đối (MAE \& MAPE)}
Bổ sung cho kiểm định tương quan, nghiên cứu tính toán các chỉ số sai số so với dữ liệu thực địa:

\begin{table}[H]
\centering
\caption{Chỉ số sai số của dữ liệu TomTom so với thực địa}
\begin{tabular}{|p{4cm}|p{3cm}|p{3cm}|p{2.5cm}|p{2cm}|}
\hline
\textbf{Chỉ số sai số} & \textbf{Đường trục chính} & \textbf{Đường nội khu} & \textbf{Giờ cao điểm} & \textbf{Tổng hợp} \\
\hline
MAE (km/h) & 1,84 & 2,31 & 2,47 & 2,11 \\
\hline
RMSE (km/h) & 2,43 & 3,07 & 3,28 & 2,81 \\
\hline
MAPE (\%) & 4,8\% & 6,1\% & 7,2\% & 5,9\% \\
\hline
Bias (km/h) --- xu hướng ước thấp & -0,32 & -0,58 & -0,71 & -0,47 \\
\hline
\end{tabular}
\end{table}

Kết quả cho thấy MAPE tổng hợp đạt 5,9\% --- nằm trong ngưỡng chấp nhận được theo tiêu chuẩn của FHWA ($< 10\%$ cho dữ liệu tốc độ đô thị). Sai số có xu hướng âm nhẹ (bias = $-0,47$ km/h), cho thấy TomTom có xu hướng ước lượng tốc độ thấp hơn thực tế một chút trong điều kiện ùn tắc --- đây là đặc điểm nhất quán với các nghiên cứu trước đó (Eastern Transportation Coalition, 2022) và không ảnh hưởng đáng kể đến giá trị phân tích.

\subsection{Luận cứ khoa học về việc sử dụng làm dữ liệu sơ cấp}
Dựa trên kết quả kiểm định định lượng ở mục 3 và hệ thống bằng chứng từ tài liệu quốc tế, phần này lập luận về tính hợp lệ của việc coi dữ liệu TomTom là dữ liệu sơ cấp trong nghiên cứu giao thông đô thị.

\subsubsection{Định nghĩa lại khái niệm dữ liệu sơ cấp trong bối cảnh Big Data}
Trong nghiên cứu khoa học xã hội và giao thông truyền thống, ``dữ liệu sơ cấp'' được định nghĩa là dữ liệu do nhà nghiên cứu trực tiếp thu thập nhằm giải quyết câu hỏi nghiên cứu cụ thể (Bryman, 2016). Tuy nhiên, trong kỷ nguyên Big Data và giao thông thông minh, khái niệm này đang được mở rộng và tái định nghĩa.

Shen và cộng sự (2020) lập luận rằng dữ liệu thụ động thu thập bởi cảm biến GPS (như FCD) hoàn toàn đáp ứng tiêu chí dữ liệu sơ cấp khi: (a) được thu thập trực tiếp từ hiện tượng nghiên cứu (phương tiện lưu thông thực tế); (b) không qua xử lý diễn giải trung gian; và (c) được thu thập theo mục đích nghiên cứu cụ thể thông qua cơ chế truy vấn API có kiểm soát. Ba điều kiện này đều được đáp ứng trong nghiên cứu hiện tại.

\subsubsection{So sánh ưu thế so với phương pháp thu thập truyền thống}

\begin{table}[H]
\centering
\caption{So sánh phương pháp thu thập dữ liệu giao thông}
\begin{tabular}{|p{3cm}|p{2.5cm}|p{2.5cm}|p{2.5cm}|p{3.5cm}|}
\hline
\textbf{Tiêu chí so sánh} & \textbf{Khảo sát thủ công} & \textbf{Vòng cảm ứng} & \textbf{Camera/CCTV} & \textbf{TomTom FCD} \\
\hline
Phạm vi không gian & Điểm cố định & Điểm cố định & Điểm/đoạn cố định & Toàn mạng lưới \\
\hline
Liên tục thời gian & Không & Có & Có & Có (24/7) \\
\hline
Chi phí triển khai & Cao & Rất cao & Cao & Thấp (API) \\
\hline
Khả năng so sánh & Thấp & Trung bình & Trung bình & Cao (chuẩn hóa) \\
\hline
Dữ liệu lịch sử & Hạn chế & Hạn chế & Hạn chế & Nhiều năm \\
\hline
Kiểm soát chất lượng & Phụ thuộc điều tra viên & Theo cơ quan quản lý & Theo nhà vận hành & Tự động, 30 giây/lần \\
\hline
\end{tabular}
\end{table}

Nguồn: Tổng hợp dựa trên Turner (2012), TomTom (2024) và khảo sát thực địa.

\subsubsection{Bằng chứng từ tài liệu học thuật quốc tế}
Việc sử dụng dữ liệu TomTom FCD làm nguồn dữ liệu chính (primary/main dataset) trong nghiên cứu đã được ghi nhận và công nhận rộng rãi trong tài liệu học thuật quốc tế:
\begin{itemize}
    \item Pozzoni và cộng sự (2023) sử dụng TomTom FCD làm dữ liệu chính để đánh giá tác động can thiệp giao thông tại Bologna (Ý), với kiểm định thống kê bằng two-tailed paired t-test, kết quả công bố trên tạp chí Sustainability (MDPI, Q2 Scimago);

    \item Eastern Transportation Coalition (2022) sử dụng dữ liệu TomTom trong báo cáo đánh giá mạng lưới giao thông liên bang Hoa Kỳ, công nhận TomTom đáp ứng tiêu chuẩn dữ liệu probe do FHWA (Federal Highway Administration) quy định;

    \item Zheng và cộng sự (2022) ứng dụng TomTom travel time data để phân tích tác động chiến tranh đến mạng lưới giao thông Ukraine, trong nghiên cứu được trích dẫn trên ResearchGate với hơn 50 lượt trích dẫn;

    \item Brederode và cộng sự (2019) trong Transportation Research Record (TRB) khẳng định TomTom FCD đáp ứng đầy đủ các tiêu chí chất lượng dữ liệu trong mô hình hóa nhu cầu giao thông.
\end{itemize}

\subsubsection{Xác nhận từ cơ quan kiểm định độc lập}
Xác nhận độc lập quan trọng nhất đến từ nghiên cứu của Đại học Michigan --- Transportation Research Institute (UMTRI), công bố năm 2013 và vẫn được trích dẫn rộng rãi đến nay. Nghiên cứu áp dụng phương pháp ``jam hunt analysis'' --- cho sáu thiết bị/ứng dụng dẫn đường chạy đồng thời qua cùng điểm ùn tắc --- và kết luận TomTom đạt độ chính xác cao nhất trong nhận diện ùn tắc với tỷ lệ 66--67\%, vượt trội so với Google Maps (52\%) và INRIX (38\%) (Belzowski, 2013).

Ở cấp độ hệ thống, Eastern Transportation Coalition (ETC) thực hiện kiểm định định kỳ và liên tục đối với dữ liệu của ba nhà cung cấp lớn (HERE, INRIX, TomTom). Các báo cáo này sử dụng WRTM (Wireless Re-identification Traffic Monitoring) làm ground truth --- phương pháp được FHWA công nhận --- và nhất quán xếp TomTom vào nhóm nhà cung cấp đáp ứng tiêu chuẩn chất lượng dữ liệu cấp quốc gia.

\subsubsection{Thảo luận hạn chế và giải pháp kiểm soát}
Để đảm bảo tính minh bạch khoa học, nghiên cứu cũng ghi nhận các hạn chế cần kiểm soát:

\paragraph{Thiên lệch phương tiện (Vehicle Type Bias)}
Dữ liệu FCD được thu thập chủ yếu từ phương tiện có trang bị thiết bị dẫn đường TomTom --- có thể không đại diện đầy đủ cho xe máy (phương tiện chiếm đa số tại đô thị Việt Nam). Giải pháp kiểm soát: phân tích riêng biệt theo loại đường và so sánh kết quả với dữ liệu đếm xe tại các điểm kiểm tra.

\paragraph{Mật độ probe thấp trên đường nhỏ}
10 đoạn đường (4,7\%) có mật độ probe dưới ngưỡng tối thiểu và đã được loại bỏ khỏi phân tích chính. Kết quả tổng hợp không bị ảnh hưởng bởi nhóm đoạn đường này.

\paragraph{Độ trễ thời gian thực}
Dữ liệu TomTom được tổng hợp và cập nhật theo chu kỳ, không phải tức thời tuyệt đối. Tuy nhiên, đối với mục tiêu nghiên cứu hành vi di chuyển và phân tích mẫu hình giao thông (không cần dự báo thời gian thực), đặc điểm này không ảnh hưởng đến giá trị sử dụng của dữ liệu.

\subsection{Kết luận}
Tổng hợp bốn nhóm bằng chứng được trình bày trong báo cáo này --- bao gồm cơ sở kỹ thuật thu thập dữ liệu, tiêu chí đánh giá chất lượng, kết quả kiểm định thống kê định lượng và xác nhận từ tài liệu học thuật quốc tế --- cho phép rút ra kết luận sau:

Dữ liệu TomTom Floating Car Data đáp ứng đầy đủ các tiêu chí độ tin cậy cần thiết để sử dụng làm dữ liệu sơ cấp trong nghiên cứu giao thông đô thị. Cụ thể: (1) mật độ probe trung bình 218,4 quan sát/giờ vượt ngưỡng khuyến nghị 7,3 lần; (2) tương quan Pearson với dữ liệu thực địa đạt $r = 0,951$ ($R^2 = 90,4\%$), phân loại ``rất mạnh'' theo thang Cohen (1988); (3) sai số MAPE = 5,9\% nằm trong ngưỡng chấp nhận; và (4) kiểm định t-test xác nhận tính nhất quán thời gian đạt 91,7\%. Các kết quả này nhất quán với tài liệu học thuật quốc tế từ UMTRI, Eastern Transportation Coalition, và nhiều nghiên cứu ứng dụng trên tạp chí peer-reviewed.

Những hạn chế đã được nhận diện và kiểm soát thông qua quy trình lọc dữ liệu nghiêm ngặt. Nghiên cứu khẳng định rằng việc sử dụng dữ liệu TomTom là lựa chọn phương pháp luận hợp lý, tiết kiệm chi phí và có cơ sở khoa học vững chắc cho các phân tích giao thông đô thị quy mô lớn.

